% =============================================================================
% Chapter: Association Rules & Apriori
% =============================================================================

\chapter{Association Rules \& Apriori}

\section{Εισαγωγή}

\keyword{Association Rule Mining} ανακαλύπτει ενδιαφέρουσες σχέσεις μεταξύ items σε μεγάλα datasets. Κλασικό παράδειγμα: \keyword{Market Basket Analysis} (ποια προϊόντα αγοράζονται μαζί).

\begin{examplebox}[Market Basket Analysis]
\begin{center}
\begin{tikzpicture}
    % Shopping basket visualization
    \node[draw=primary, fill=box-yellow, rounded corners=5pt, minimum width=8cm, minimum height=2cm] (basket) {};
    \node[anchor=north, font=\bfseries] at (basket.north) {Transaction};

    % Items
    \node[draw=success, fill=success!20, rounded corners=3pt, minimum width=1.5cm] at (-2.5,0) {Bread};
    \node[draw=secondary, fill=secondary!20, rounded corners=3pt, minimum width=1.5cm] at (-0.5,0) {Butter};
    \node[draw=info, fill=info!20, rounded corners=3pt, minimum width=1.5cm] at (1.5,0) {Milk};

    % Arrow and rule
    \draw[->, >=Stealth, thick] (5,0) -- (7,0);
    \node[anchor=west, align=left] at (7.2,0) {
        \textbf{Rule}: \{Bread, Butter\} $\to$ \{Milk\}\\
        {\small "Αν αγοράσεις ψωμί και βούτυρο,}\\
        {\small πιθανώς θα αγοράσεις και γάλα"}
    };
\end{tikzpicture}
\end{center}
\end{examplebox}

% =============================================================================
\section{Βασικές Έννοιες \& Ορισμοί}

\subsection{Itemset}

\begin{definitionbox}[Itemset]
\textbf{Itemset}: Σύνολο από items (π.χ. \{Bread, Milk, Butter\})

\textbf{k-itemset}: Itemset με k items
\begin{itemize}
    \item 1-itemset: \{Bread\}
    \item 2-itemset: \{Bread, Milk\}
    \item 3-itemset: \{Bread, Milk, Butter\}
\end{itemize}
\end{definitionbox}

\subsection{Transaction Database}

\begin{center}
\begin{tabular}{cl}
\toprule
\textbf{TID} & \textbf{Items} \\
\midrule
T1 & Bread, Milk \\
T2 & Bread, Butter, Eggs \\
T3 & Milk, Butter, Eggs \\
T4 & Bread, Milk, Butter, Eggs \\
T5 & Bread, Milk, Butter \\
\bottomrule
\end{tabular}

\vspace{0.3cm}
Total Transactions: 5
\end{center}

% =============================================================================
\section{Μέτρα Αξιολόγησης}

\subsection{Support}

\begin{formulabox}[Support]
\begin{equation}
\text{Support}(X) = \frac{|\{T : X \subseteq T\}|}{|T|} = \frac{\text{Transactions που περιέχουν } X}{\text{Συνολικές Transactions}}
\end{equation}

"Πόσο συχνά εμφανίζεται το itemset X στις transactions"
\end{formulabox}

\begin{examplebox}[Υπολογισμός Support]
\begin{align*}
\text{Support}(\{Bread\}) &= 4/5 = 0.80 \quad \text{(T1, T2, T4, T5)} \\
\text{Support}(\{Milk\}) &= 4/5 = 0.80 \quad \text{(T1, T3, T4, T5)} \\
\text{Support}(\{Bread, Milk\}) &= 3/5 = 0.60 \quad \text{(T1, T4, T5)} \\
\text{Support}(\{Bread, Butter\}) &= 3/5 = 0.60 \quad \text{(T2, T4, T5)}
\end{align*}
\end{examplebox}

\subsection{Confidence}

\begin{formulabox}[Confidence]
\begin{equation}
\text{Confidence}(X \to Y) = \frac{\text{Support}(X \cup Y)}{\text{Support}(X)}
\end{equation}

"Αν έχω X, πόσο πιθανό να έχω και Y;"
\end{formulabox}

\begin{examplebox}[Υπολογισμός Confidence]
\begin{align*}
\text{Confidence}(\{Bread\} \to \{Milk\}) &= \frac{\text{Support}(\{Bread, Milk\})}{\text{Support}(\{Bread\})} \\
&= \frac{0.60}{0.80} = 0.75 = 75\%
\end{align*}

"75\% των transactions με Bread περιέχουν και Milk"
\end{examplebox}

\subsection{Lift}

\begin{formulabox}[Lift]
\begin{equation}
\text{Lift}(X \to Y) = \frac{\text{Confidence}(X \to Y)}{\text{Support}(Y)} = \frac{\text{Support}(X \cup Y)}{\text{Support}(X) \times \text{Support}(Y)}
\end{equation}
\end{formulabox}

\begin{center}
\begin{tikzpicture}
    % Lift interpretation visualization
    \node[draw=accent, fill=accent!20, rounded corners=5pt, minimum width=3.5cm, minimum height=1.8cm] (neg) at (0,0) {};
    \node[anchor=center, align=center] at (neg.center) {\textbf{Lift $<$ 1}\\Αρνητική\\συσχέτιση};

    \node[draw=warning, fill=warning!20, rounded corners=5pt, minimum width=3.5cm, minimum height=1.8cm] (ind) at (5,0) {};
    \node[anchor=center, align=center] at (ind.center) {\textbf{Lift $=$ 1}\\Ανεξάρτητα\\(τυχαίο)};

    \node[draw=success, fill=success!20, rounded corners=5pt, minimum width=3.5cm, minimum height=1.8cm] (pos) at (10,0) {};
    \node[anchor=center, align=center] at (pos.center) {\textbf{Lift $>$ 1}\\Θετική\\συσχέτιση};

    % Scale
    \draw[thick, ->, >=Stealth] (-2,-1.8) -- (12,-1.8);
    \node at (0,-2.2) {0};
    \node at (5,-2.2) {1};
    \node at (10,-2.2) {$\infty$};
    \node[anchor=north] at (5,-2.5) {\small Lift Value};
\end{tikzpicture}
\end{center}

\subsection{Σύνοψη Μέτρων}

\begin{center}
\begin{tikzpicture}
    \matrix (m) [
        matrix of nodes,
        nodes={
            draw=primary,
            minimum width=4cm,
            minimum height=1cm,
            anchor=center,
            font=\small
        },
        column 1/.style={nodes={fill=ntua-blue!20, font=\bfseries}},
        row 1/.style={nodes={fill=secondary!30, font=\bfseries}},
        column sep=-\pgflinewidth,
        row sep=-\pgflinewidth
    ] {
        Μέτρο & Τύπος & Ερμηνεία \\
        Support(X) & $|T \text{ with } X| / |T|$ & Συχνότητα εμφάνισης \\
        Confidence(X$\to$Y) & $\frac{\text{Sup}(X \cup Y)}{\text{Sup}(X)}$ & Conditional probability \\
        Lift(X$\to$Y) & $\frac{\text{Conf}(X \to Y)}{\text{Sup}(Y)}$ & Correlation measure \\
    };
\end{tikzpicture}
\end{center}

% =============================================================================
\section{Apriori Algorithm}

\subsection{Βασική Αρχή: Downward Closure (Apriori Principle)}

\begin{warningbox}[Apriori Principle - ΑΠΟΜΝΗΜΟΝΕΥΣΕ!]
\begin{center}
\begin{tikzpicture}
    % Main principle box
    \node[draw=accent, fill=box-red, rounded corners=8pt, minimum width=12cm, minimum height=3cm] (principle) {};

    % Text
    \node[anchor=north, align=center, font=\large] at ([yshift=-0.3cm]principle.north) {
        \textbf{"Αν ένα itemset είναι INFREQUENT,}\\
        \textbf{όλα τα SUPERSETS του είναι επίσης infrequent"}
    };

    \node[anchor=south, align=center, font=\small] at ([yshift=0.3cm]principle.south) {
        ΑΝΤΙΣΤΡΟΦΑ: "Κάθε subset ενός frequent itemset είναι frequent"
    };
\end{tikzpicture}
\end{center}

\textbf{Γιατί είναι χρήσιμο}: Αποκλείουμε candidates χωρίς να μετρήσουμε support!
\end{warningbox}

\begin{center}
\begin{tikzpicture}[
    level distance=1.5cm,
    sibling distance=3cm,
    edge from parent/.style={draw, ->, >=Stealth}
]
    % Lattice visualization
    \node[draw=accent, fill=accent!30, circle, minimum size=1cm] (ab) {\{A,B\}}
        [grow=down]
        child {
            node[draw=accent, fill=accent!20, circle, minimum size=1cm] (abc) {\{A,B,C\}}
            child {
                node[draw=accent, fill=accent!10, circle, minimum size=0.8cm] (abcd) {\small\{A,B,C,D\}}
            }
        };

    % Labels
    \node[anchor=west, font=\small, text=accent] at (2,0) {Infrequent};
    \node[anchor=west, font=\small, text=accent] at (2,-1.5) {$\Rightarrow$ Also Infrequent};
    \node[anchor=west, font=\small, text=accent] at (2,-3) {$\Rightarrow$ Also Infrequent};

    % Support inequality
    \node[anchor=north, align=center] at (0,-4.5) {
        $\text{Support}(\{A,B,C,D\}) \leq \text{Support}(\{A,B,C\}) \leq \text{Support}(\{A,B\})$
    };
\end{tikzpicture}
\end{center}

\subsection{Αλγόριθμος}

\begin{algorithmbox}[Apriori Algorithm]
\begin{lstlisting}[language=pseudo]
function APRIORI(D, min_sup):
    // D = transaction database
    // min_sup = minimum support threshold

    C1 = all 1-itemsets
    L1 = {items in C1 with support >= min_sup}

    k = 2
    while L(k-1) is not empty:
        // Generate candidates
        Ck = generate_candidates(L(k-1))

        // Prune using Apriori principle
        Ck = prune(Ck)  // remove if any (k-1)-subset not in L(k-1)

        // Count support
        for each transaction t in D:
            for each candidate c in Ck:
                if c is subset of t:
                    count[c]++

        // Filter frequent
        Lk = {c in Ck : support(c) >= min_sup}
        k = k + 1

    return L1 union L2 union ... union L(k-1)
\end{lstlisting}
\end{algorithmbox}

\subsection{Apriori Flowchart}

\begin{center}
\begin{tikzpicture}[node distance=1.5cm]
    % Start
    \node[startstop] (start) {Start};

    % Scan C1
    \node[process, below=of start] (c1) {$C_1$ = all 1-itemsets};

    % L1
    \node[process, below=of c1] (l1) {$L_1$ = frequent 1-itemsets};

    % k=2
    \node[process, below=of l1] (k2) {k = 2};

    % Check empty
    \node[decision, below=of k2, aspect=2] (check) {$L_{k-1} = \emptyset$?};

    % Generate Ck
    \node[process, right=3cm of check] (gen) {Generate $C_k$ from $L_{k-1}$};

    % Prune
    \node[process, below=of gen] (prune) {Prune $C_k$};

    % Count support
    \node[process, below=of prune] (count) {Count support of $C_k$};

    % Lk
    \node[process, below=of count] (lk) {$L_k$ = frequent k-itemsets};

    % k++
    \node[process, left=3cm of lk] (inc) {k = k + 1};

    % End
    \node[startstop, below=of check] (end) {Return $\bigcup L_k$};

    % Arrows
    \draw[arrow] (start) -- (c1);
    \draw[arrow] (c1) -- (l1);
    \draw[arrow] (l1) -- (k2);
    \draw[arrow] (k2) -- (check);
    \draw[arrow] (check) -- node[above] {No} (gen);
    \draw[arrow] (gen) -- (prune);
    \draw[arrow] (prune) -- (count);
    \draw[arrow] (count) -- (lk);
    \draw[arrow] (lk) -- (inc);
    \draw[arrow] (inc) |- (check);
    \draw[arrow] (check) -- node[right] {Yes} (end);
\end{tikzpicture}
\end{center}

% =============================================================================
\section{Worked Example (ΤΥΠΙΚΟ ΘΕΜΑ ΕΞΕΤΑΣΕΩΝ!)}

\begin{examplebox}[Apriori - Πλήρες Παράδειγμα]
\textbf{Transaction Database:}
\begin{center}
\begin{tabular}{cl}
\toprule
\textbf{TID} & \textbf{Items} \\
\midrule
T1 & A, C, D \\
T2 & B, C, E \\
T3 & A, B, C, E \\
T4 & B, E \\
T5 & A, B, C, E \\
\bottomrule
\end{tabular}

\vspace{0.3cm}
$\text{min\_sup} = 2/5 = 0.40$ (τουλάχιστον 2 transactions)
\end{center}
\end{examplebox}

\subsection{Βήμα 1: $C_1 \to L_1$}

\begin{center}
\begin{tikzpicture}
    % C1 box
    \node[draw=secondary, fill=box-blue, rounded corners=5pt, minimum width=5cm, minimum height=3cm] (c1) at (0,0) {};
    \node[anchor=north, font=\bfseries] at (c1.north) {$C_1$ (Candidates)};
    \node[anchor=center, align=left, font=\small] at (c1.center) {
        \{A\}: 3 $\to$ 0.60 $\geq$ 0.40 \textcolor{success}{\checkmark}\\
        \{B\}: 4 $\to$ 0.80 $\geq$ 0.40 \textcolor{success}{\checkmark}\\
        \{C\}: 4 $\to$ 0.80 $\geq$ 0.40 \textcolor{success}{\checkmark}\\
        \{D\}: 1 $\to$ 0.20 $<$ 0.40 \textcolor{accent}{\texttimes}\\
        \{E\}: 4 $\to$ 0.80 $\geq$ 0.40 \textcolor{success}{\checkmark}
    };

    % Arrow
    \draw[->, >=Stealth, thick] (3,0) -- (5,0) node[midway, above] {Filter};

    % L1 box
    \node[draw=success, fill=box-green, rounded corners=5pt, minimum width=4cm, minimum height=2cm] (l1) at (7.5,0) {};
    \node[anchor=north, font=\bfseries] at (l1.north) {$L_1$ (Frequent)};
    \node[anchor=center, align=center] at (l1.center) {
        \{A\}, \{B\}, \{C\}, \{E\}
    };
\end{tikzpicture}
\end{center}

\subsection{Βήμα 2: Generate $C_2$ from $L_1$}

\begin{center}
\begin{tikzpicture}
    % C2 box
    \node[draw=secondary, fill=box-blue, rounded corners=5pt, minimum width=6cm, minimum height=4cm] (c2) at (0,0) {};
    \node[anchor=north, font=\bfseries] at (c2.north) {$C_2$ (2-combinations of $L_1$)};
    \node[anchor=center, align=left, font=\small] at (c2.center) {
        \{A,B\}: 2 $\to$ 0.40 \textcolor{success}{\checkmark}\\
        \{A,C\}: 3 $\to$ 0.60 \textcolor{success}{\checkmark}\\
        \{A,E\}: 2 $\to$ 0.40 \textcolor{success}{\checkmark}\\
        \{B,C\}: 3 $\to$ 0.60 \textcolor{success}{\checkmark}\\
        \{B,E\}: 4 $\to$ 0.80 \textcolor{success}{\checkmark}\\
        \{C,E\}: 3 $\to$ 0.60 \textcolor{success}{\checkmark}
    };

    % Arrow
    \draw[->, >=Stealth, thick] (4,0) -- (6,0) node[midway, above] {All pass!};

    % L2 box
    \node[draw=success, fill=box-green, rounded corners=5pt, minimum width=5.5cm, minimum height=2.5cm] (l2) at (9,0) {};
    \node[anchor=north, font=\bfseries] at (l2.north) {$L_2$ (Frequent)};
    \node[anchor=center, align=center, font=\small] at (l2.center) {
        \{A,B\}, \{A,C\}, \{A,E\}\\
        \{B,C\}, \{B,E\}, \{C,E\}
    };
\end{tikzpicture}
\end{center}

\subsection{Βήμα 3: Generate $C_3$ from $L_2$}

\textbf{Join Step:}
\begin{align*}
\{A,B\} \cup \{A,C\} &= \{A,B,C\} \\
\{A,B\} \cup \{A,E\} &= \{A,B,E\} \\
\{A,C\} \cup \{A,E\} &= \{A,C,E\} \\
\{B,C\} \cup \{B,E\} &= \{B,C,E\}
\end{align*}

\textbf{Prune Step:}
\begin{center}
\begin{tabular}{llll}
\toprule
\textbf{Candidate} & \textbf{Subsets in $L_2$?} & \textbf{Status} \\
\midrule
\{A,B,C\} & \{A,B\}\checkmark, \{A,C\}\checkmark, \{B,C\}\checkmark & \textcolor{success}{KEEP} \\
\{A,B,E\} & \{A,B\}\checkmark, \{A,E\}\checkmark, \{B,E\}\checkmark & \textcolor{success}{KEEP} \\
\{A,C,E\} & \{A,C\}\checkmark, \{A,E\}\checkmark, \{C,E\}\checkmark & \textcolor{success}{KEEP} \\
\{B,C,E\} & \{B,C\}\checkmark, \{B,E\}\checkmark, \{C,E\}\checkmark & \textcolor{success}{KEEP} \\
\bottomrule
\end{tabular}
\end{center}

\textbf{Count Support:}
\[
L_3 = \{\{A,B,C\}, \{A,B,E\}, \{A,C,E\}, \{B,C,E\}\}
\]

\subsection{Βήμα 4: Generate $C_4$ from $L_3$}

Join: $\{A,B,C\} \cup \{A,B,E\} = \{A,B,C,E\}$

Prune: All 3-subsets in $L_3$? \checkmark

Count: \{A,B,C,E\} appears in T3, T5 $\to$ support = 0.40 $\geq$ 0.40 \checkmark

\[
L_4 = \{\{A,B,C,E\}\}
\]

\subsection{Τελικό Αποτέλεσμα}

\begin{center}
\begin{tikzpicture}
    \node[draw=ntua-blue, fill=ntua-blue!10, rounded corners=10pt, minimum width=12cm, minimum height=5cm] (result) {};
    \node[anchor=north, font=\Large\bfseries\color{ntua-blue}] at (result.north) {Frequent Itemsets};

    \node[anchor=center, align=left] at (result.center) {
        \textbf{1-itemsets}: \{A\}, \{B\}, \{C\}, \{E\}\\[0.5em]
        \textbf{2-itemsets}: \{A,B\}, \{A,C\}, \{A,E\}, \{B,C\}, \{B,E\}, \{C,E\}\\[0.5em]
        \textbf{3-itemsets}: \{A,B,C\}, \{A,B,E\}, \{A,C,E\}, \{B,C,E\}\\[0.5em]
        \textbf{4-itemsets}: \{A,B,C,E\}
    };
\end{tikzpicture}
\end{center}

% =============================================================================
\section{Rule Generation}

\begin{formulabox}[Από Frequent Itemset σε Rules]
Για κάθε frequent itemset $L$, για κάθε non-empty subset $S$:
\begin{itemize}
    \item Rule: $(L - S) \to S$
    \item $\text{Confidence} = \text{Support}(L) / \text{Support}(L - S)$
\end{itemize}
\end{formulabox}

\begin{examplebox}[Rule Generation]
Frequent itemset: \{A, B, C\} με support = 0.40

\begin{center}
\begin{tabular}{lcc}
\toprule
\textbf{Rule} & \textbf{Confidence} & \textbf{Pass?} (min\_conf=0.60) \\
\midrule
\{A, B\} $\to$ \{C\} & $0.40/0.40 = 1.00$ & \textcolor{success}{\checkmark} \\
\{A, C\} $\to$ \{B\} & $0.40/0.60 = 0.67$ & \textcolor{success}{\checkmark} \\
\{B, C\} $\to$ \{A\} & $0.40/0.60 = 0.67$ & \textcolor{success}{\checkmark} \\
\{A\} $\to$ \{B, C\} & $0.40/0.60 = 0.67$ & \textcolor{success}{\checkmark} \\
\{B\} $\to$ \{A, C\} & $0.40/0.80 = 0.50$ & \textcolor{accent}{\texttimes} \\
\{C\} $\to$ \{A, B\} & $0.40/0.80 = 0.50$ & \textcolor{accent}{\texttimes} \\
\bottomrule
\end{tabular}
\end{center}
\end{examplebox}

% =============================================================================
\section{FP-Growth (Overview)}

\begin{center}
\begin{tikzpicture}
    % Apriori box
    \node[draw=accent, fill=accent!20, rounded corners=10pt, minimum width=5.5cm, minimum height=4cm] (apriori) at (0,0) {};
    \node[anchor=north, font=\bfseries\color{accent}] at (apriori.north) {Apriori};
    \node[anchor=center, align=left, font=\small] at (apriori.center) {
        \textbf{Μειονεκτήματα:}\\
        • Πολλαπλά scans βάσης\\
        • Μεγάλος αριθμός candidates\\
        • Πολλοί support counts
    };

    % Arrow
    \draw[->, >=Stealth, thick, dashed] (3.5,0) -- (5.5,0) node[midway, above] {vs};

    % FP-Growth box
    \node[draw=success, fill=success!20, rounded corners=10pt, minimum width=5.5cm, minimum height=4cm] (fpg) at (9,0) {};
    \node[anchor=north, font=\bfseries\color{success}] at (fpg.north) {FP-Growth};
    \node[anchor=center, align=left, font=\small] at (fpg.center) {
        \textbf{Πλεονεκτήματα:}\\
        • Μόνο 2 database scans\\
        • Χωρίς candidate generation\\
        • FP-tree compression
    };
\end{tikzpicture}
\end{center}

% =============================================================================
\section{Συχνά Λάθη}

\begin{warningbox}[Κοινά Λάθη στις Εξετάσεις]
\begin{enumerate}
    \item \textbf{Support Calculation}:
    \begin{itemize}
        \item[$\times$] $\text{Support}(\{A,B\}) = \text{Support}(\{A\}) \times \text{Support}(\{B\})$
        \item[\checkmark] $\text{Support}(\{A,B\}) = |T \text{ containing BOTH A and B}| / |T|$
    \end{itemize}

    \item \textbf{Confidence Direction}:
    \begin{itemize}
        \item[$\times$] $\text{Confidence}(A \to B) = \text{Confidence}(B \to A)$
        \item[\checkmark] Συνήθως διαφορετικά! Εξαρτάται από Support
    \end{itemize}

    \item \textbf{Lift Interpretation}:
    \begin{itemize}
        \item Υψηλό confidence ΔΕΝ σημαίνει υψηλό lift!
        \item Lift = 1 σημαίνει ανεξαρτησία
    \end{itemize}
\end{enumerate}
\end{warningbox}

% =============================================================================
\section{Σύνοψη - Key Points}

\begin{center}
\begin{tikzpicture}
    % Main summary box
    \node[draw=ntua-blue, fill=ntua-blue!10, rounded corners=5pt, minimum width=14cm, minimum height=9cm] (main) at (0,0) {};

    % Formulas section
    \node[anchor=north west, font=\bfseries\Large\color{ntua-blue}] at ([xshift=5pt, yshift=-5pt]main.north west) {Τύποι (ΑΠΟΜΝΗΜΟΝΕΥΣΕ!)};

    \node[anchor=north west, align=left] at ([xshift=10pt, yshift=-35pt]main.north west) {
        \begin{tabular}{ll}
            \textbf{Support(X)} & $= |T \text{ containing } X| / |T|$ \\[8pt]
            \textbf{Confidence(X$\to$Y)} & $= \text{Support}(X \cup Y) / \text{Support}(X)$ \\[8pt]
            \textbf{Lift(X$\to$Y)} & $= \text{Conf}(X \to Y) / \text{Sup}(Y)$ \\[8pt]
             & $= \text{Sup}(X \cup Y) / (\text{Sup}(X) \times \text{Sup}(Y))$
        \end{tabular}
    };

    % Apriori steps
    \node[anchor=west, align=left] at ([xshift=10pt, yshift=-60pt]main.west) {
        \textbf{Apriori Steps:}\\
        1. Find frequent 1-itemsets ($L_1$)\\
        2. Generate candidates ($C_k$) from $L_{k-1}$\\
        3. Prune candidates (Apriori principle)\\
        4. Count support, keep frequent ($L_k$)\\
        5. Repeat until no more frequent itemsets
    };

    % Apriori Principle
    \node[draw=accent, fill=box-red, rounded corners=5pt, minimum width=6cm, minimum height=1.5cm]
        at ([xshift=-3.5cm, yshift=1.5cm]main.south) {};
    \node[align=center, font=\small] at ([xshift=-3.5cm, yshift=1.5cm]main.south) {
        \textbf{Apriori Principle:}\\
        "Subsets of frequent itemsets are frequent"
    };

    % Lift box
    \node[draw=success, fill=box-green, rounded corners=5pt, minimum width=4.5cm, minimum height=1.5cm]
        at ([xshift=3.5cm, yshift=1.5cm]main.south) {};
    \node[align=center, font=\small] at ([xshift=3.5cm, yshift=1.5cm]main.south) {
        \textbf{Lift:} 1 = Independent\\
        $>$1 = Positive, $<$1 = Negative
    };
\end{tikzpicture}
\end{center}
